% !TeX root = ../main.tex

% ==================================================
% 2.5. 메서드와 self — 인스턴스 메서드는 어떻게 호출되는가?
% ==================================================

\subsection{메서드와 \zsh{self} - 인스턴스 메서드는 어떻게 호출되는가?}

2.4에서는 인스턴스 속성과 네임스페이스를 학습했다.

특히 \zsh{self.name = name}이 지역 변수 \zsh{name} 의 값을 인스턴스 속성에 연결하는 동작이라는 점을 확인했다.

이번 2.5에서는 클래스에 정의된 함수가 인스턴스를 통해 호출될 때 어떤 일이 발생하는지 살펴볼 것이다.

이번 학습의 출발점은 다음 코드다.

\begin{dhkcode}
class Person:
    def introduce(self):
        print("Hello")

p = Person()
p.introduce()
\end{dhkcode}

여기서 가장 중요한 질문은 하나다.

\zsh{introduce()}에는 \zsh{self}라는 매개변수가 있는데, 왜 \zsh{p.introduce()}에서는 인수를 전달하지 않아도 되는가?

이 질문에 답하려면 함수 객체, 속성 탐색, 바운드 메서드(Bound Method), 인수 바인딩을 연결해서 이해해야 한다.

\subsubsection{개념 정의 - 함수와 메서드는 무엇이 다른가?}

먼저 일반 함수를 살펴보자.

\begin{dhkcode}
def introduce():
    print("Hello")
\end{dhkcode}

이 코드를 실행하면 Python은 함수 객체를 생성하고 \zsh{introduce}라는 이름에 바인딩한다.

\begin{dhkcode}
introduce()
\end{dhkcode}

이 표현식은 \zsh{introduce}라는 이름이 참조하는 함수 객체를 호출하는 것이다.

그렇다면 다음 코드는 어떨까?

\begin{dhkcode}
class Person:
    def introduce(self):
        print("Hello")
\end{dhkcode}

여기서도 \zsh{def}가 실행되면 함수 객체가 생성된다.

다만 그 함수 객체가 모듈의 전역 이름이 아니라 클래스의 네임스페이스에 \zsh{introduce}라는 이름으로 저장된다는 차이가 있다.

\begin{dhkcode}
print(Person.__dict__["introduce"])
\end{dhkcode}

출력 예시:

\begin{dhkzsh}
<function Person.introduce at 0x...>
\end{dhkzsh}

즉, \zsh{Person} 클래스에 정의된 \zsh{introduce}도 본래는 함수 객체다.

그런데 이 함수를 인스턴스를 통해 접근하면 Python의 메서드 바인딩 규칙이 적용된다.

\begin{dhktable}
|표현식|	얻는 객체|
|:---|:---|
|Person.__dict__["introduce"]|	클래스에 저장된 원래 함수 객체|
|Person.introduce|	이 예제에서는 같은 함수 객체|
|p.introduce|	p에 바인딩된 메서드 객체|
|p.introduce()|	바운드 메서드를 호출한 결과|
\end{dhktable}

여기서 중요한 구분은 클래스에 함수가 저장되어 있다는 사실과, 인스턴스를 통해 그 함수에 접근할 때 메서드 바인딩이 일어난다는 사실이다.

\subsubsection{메서드는 왜 필요한가?}

다음과 같이 사람의 이름을 출력하는 일반 함수를 생각해 보자.

\begin{dhkcode}
def introduce(person):
    print(f"Hello, {person.name}")

class Person:
    pass

p = Person()
p.name = "Kim"

introduce(p)
\end{dhkcode}

실행 결과:

\begin{dhkzsh}
Hello, Kim
\end{dhkzsh}

여기서는 \zsh{introduce()}가 어떤 사람의 정보를 사용할지 알 수 있도록 \zsh{p}를 인수로 전달했다.

그런데 클래스에 함수를 정의하면 다음과 같이 작성할 수 있다.

\begin{dhkcode}
class Person:
    def introduce(self):
        print(f"Hello, {self.name}")

p = Person()
p.name = "Kim"

p.introduce()
\end{dhkcode}

실행 결과:

\begin{dhkzsh}
Hello, Kim
\end{dhkzsh}

이 코드는 \zsh{p}라는 객체에 연결된 동작을 호출한다는 의미를 명확하게 표현한다.

객체지향 프로그래밍에서는 이렇게 객체의 상태와 그 상태를 사용하는 동작을 하나의 인터페이스로 구성할 수 있다.

다만 Python에서는 클래스에 정의된 모든 함수가 반드시 인스턴스의 상태를 사용해야 하는 것은 아니다.

인스턴스 상태가 필요 없는 경우에는 정적 메서드 등을 사용할 수도 있으며 이 내용은 2.11에서 다룬다.

\subsubsection{\zsh{self}는 무엇인가?}

다음 코드를 다시 보자.

\begin{dhkcode}
class Person:
    def introduce(self):
        print(self.name)

p = Person()
p.name = "Kim"

p.introduce
\end{dhkcode}

여기서 \zsh{self}는 메서드 호출의 대상이 되는 인스턴스를 전달받는 매개변수다.

\zsh{self} 자체는 Python의 예약어가 아니다.

예를 들어 다음 코드도 문법적으로 가능하다.

\begin{dhkcode}
class Person:
    def introduce(this):
        print(this.name)

p = Person()
p.name = "Kim"

p.introduce()
\end{dhkcode}

결과는 동일하게 \zsh{Kim}이 출력된다.

하지만 Python에서는 인스턴스 메서드의 첫 번째 매개변수 이름으로 \zsh{self}를 사용하는 것이 관례이다.

따라서 특별한 이유가 없다면 반드시 이 관례를 따르는 것이 좋다.

\begin{dhknote}
\begin{dhksmry}[centered]
self는 클래스 객체를 의미하는가?
\end{dhksmry}
\end{dhknote}

아니다.

\begin{dhkcode}
class Person:
    def introduce(self):
        print(self)
        print(type(self))

p = Person()
p.introduce()
\end{dhkcode}

출력 예시:

\begin{dhkzsh}
<__main__.Person object at 0x...>
<class '__main__.Person'>
\end{dhkzsh}

\zsh{self}는 \zsh{Person} 클래스 객체가 아니라 \zsh{p}가 참조하는 인스턴스 객체다.

다음 코드로 동일성을 확인할 수도 있다.

\begin{dhkcode}
class Person:
    def check(self, other):
        print(self is other)


p = Person()
p.check(p)
\end{dhkcode}

실행 결과:

\begin{dhkzsh}
True
\end{dhkzsh}

이 호출에서 \zsh{self}와 \zsh{other}는 모두 동일한 인스턴스 객체를 참조한다.

다만 두 매개변수는 별개의 지역 이름이다.

\subsubsection{\zsh{p.introduce()}는 어떻게 실행되는가?}

이제 이번 학습의 핵심으로 들어가 보자.

\begin{dhkcode}
class Person:
    def introduce(self):
        print("Hello")

p = Person()

p.introduce()
\end{dhkcode}

Python은 \zsh{p.introduce()}를 처리할 때 다음 과정을 거친다.

\begin{dhkgraph}[uniform]
node i1 (1,1) :: p.introduce() // 속성 접근 및 호출 표현식
node i2 (2,1) :: p.introduce // 인스턴스를 통한 속성 접근
node i3 (3,1) :: 함수 객체 탐색 // Person 클래스에서 introduce 발견
node i4 (4,1) [accent] :: 바운드 메서드 생성 // 함수 객체와 p 인스턴스를 연결
node i5 (5,1) :: 메서드 호출 // self에 p를 전달해 함수 본문 실행

edge i1 -> i2
edge i2 -> i3
edge i3 -> i4
edge i4 -> i5

caption :: 일반적인 인스턴스 메서드 호출 과정을 단순화한 그림
\end{dhkgraph}

이 과정에서 중요한 부분은 세 번째와 네 번째 단계다.

Python은 클래스에서 \zsh{introduce} 함수 객체를 찾은 뒤, 인스턴스 \zsh{p}와 연결된 바운드 메서드를 제공한다.

그 결과 \zsh{p.introduce()}를 호출하면 \zsh{p}가 첫 번째 인수로 전달되는 효과가 발생한다.

따라서 다음 두 호출은 이 예제에서 동일한 동작을 수행한다.

\begin{dhkcode}
p.introduce()

Person.introduce(p)
\end{dhkcode}

두 번째 코드는 클래스에 저장된 함수 객체를 직접 호출하면서 \zsh{p}를 명시적으로 전달한 것이다.

첫 번째 코드는 Python이 인스턴스와 함수를 연결한 바운드 메서드를 통해 호출한 것이다.

직접 확인해 보자.

\begin{dhkcode}
class Person:
    def introduce(self):
        print("Hello")

p = Person()
print(Person.introduce)
print(p.introduce)
\end{dhkcode}

출력 예시:

\begin{dhkzsh}
<function Person.introduce at 0x...>
<bound method Person.introduce of <__main__.Person object at 0x...>>
\end{dhkzsh}

여기서 두 번째 출력에 주목한다.

\begin{dhkzsh}
bound method
\end{dhkzsh}

이 것이 바로 바운드 메서드다.

\subsubsection{바운드 메서드의 내부 구조 - \zsh{__self__}와 \zsh{__func__}}

바운드 메서드는 단순히 함수 이름을 다르게 표현한 것이 아니다.

원래 함수 객체와 연결된 인스턴스 객체를 함께 참조하는 별도의 메서드 객체다.

다음 코드를 보자.

\begin{dhkcode}
class Person:
    def introduce(self):
        print(self.name)

p = Person()
p.name = "Kim"

method = p.introduce
\end{dhkcode}

이제 \zsh{method}가 무엇을 참조하는지 확인해 보자.

\begin{dhkcode}
print(method.__self__ is p)
print(method.__func__ is Person.introduce)
\end{dhkcode}

실행 결과:

\begin{dhkzsh}
True
True
\end{dhkzsh}

각 속성의 의미는 다음과 같다.

\begin{dhktable}
|속성|의미|
|:---|:---|
|method.__self__|메서드에 연결된 인스턴스 객체|
|method.__func__|실제 호출할 원래 함수 객체|
\end{dhktable}

이를 그림으로 나타내면 다음과 같다.

\begin{dhkgraph}
node m (1,2) [accent] :: method//바운드 메서드
node s (2,1) :: __self__//p 인스턴스//name = "kim"
node f (2,3) :: __func__//Person.introduce//함수 객체

split m -> s, f

caption::바운드 메서드는 호출 대상 인스턴스와 원래 함수 객체를 참조한다.
\end{dhkgraph}

따라서 다음 코드를 실행하면

\begin{dhkcode}
method()
\end{dhkcode}

개념적으로 다음과 같은 호출이 이루어진다.

\begin{dhkcode}
method.__func__(method.__self__)
\end{dhkcode}

즉,

\begin{dhkcode}
Person.introduce(p)
\end{dhkcode}

와 같은 효과를 갖는 것이다.

여기서 이전에 배운 함수 객체와 이름의 바인딩 개념이 다시 등장한다.

\zsh{method}는 바운드 메서드 객체를 참조하는 이름일 뿐이다.

\begin{dhknote}
\begin{dhksmry}[centered]
바운드 메서드는 변수에 저장할 수도 있다.
\end{dhksmry}
\end{dhknote}

\begin{dhkcode}
method = p.introduce

method()
\end{dhkcode}

실행 결과:

\begin{dhkzsh}
Kim
\end{dhkzsh}

이것이 가능한 이유는 바운드 메서드 역시 Callable 객체이기 때문이다.

또한 \zsh{method}가 인스턴스 \zsh{p}를 참조하고 있으므로, 이후에 \zsh{p}라는 이름을 다른 객체에 재바인딩하더라도 기존 바운드 메서드는 원래 인스턴스에 연결되어 있다.

\begin{dhkcode}
original = p
method = p.introduce

p = Person()
p.name = "Lee"

method()
\end{dhkcode}

실행 결과:

\begin{dhkzsh}
Kim
\end{dhkzsh}

\zsh{method}는 새로운 \zsh{p}가 아니라 원래 인스턴스를 참조하기 때문이다.

\subsubsection{메서드 바인딩은 왜 발생하는가? - 디스크립터 프로토콜}

여기서 조금 더 깊게 들어가 보자.

앞에서 Python은 인스턴스를 통해 함수에 접근할 때 바운드 메서드를 제공한다고 설명하였다.

그런데 Python은 어떻게 클래스에 저장된 함수 객체를 바운드 메서드로 변환할까?

그 핵심에는 \zsh{디스크립터 프로토콜(Descriptor Protocol)}이 있다.

Python의 일반적인 함수 객체는 \zsh{__get__()} 메서드를 제공한다.

클래스에 저장된 함수 객체에 인스턴스를 통해 접근하면, 속성 탐색 과정에서 이 \zsh{__get__()}이 활용되어 바운드 메서드가 만들어진다.

다음 코드를 살펴보자.

\begin{dhkcode}
class Person:
    def introduce(self):
        print("Hello")

p = Person()

func = Person.__dict__["introduce"]
\end{dhkcode}

여기서 \zsh{func}는 클래스에 저장된 원래 함수 객체다.

다음 코드를 실행해 보자.

\begin{dhkcode}
method = func.__get__(p, Person)

print(method)
\end{dhkcode}

출력 예시:

\begin{dhkzsh}
<bound method Person.introduce of <__main__.Person object at 0x...>>
\end{dhkzsh}

이제 다음 두 표현식을 비교해 보자.

\begin{dhkcode}
p.introduce

func.__get__(p, Person)
\end{dhkcode}

이 예제에서는 둘 다 \zsh{p}에 연결된 바운드 메서드를 얻는 효과가 있다.

실제로 확인해 보자.

\begin{dhkcode}
method = func.__get__(p, Person)

print(method.__self__ is p)
print(method.__func__ is func)
\end{dhkcode}

실행 결과:

\begin{dhkzsh}
True
True
\end{dhkzsh}

따라서 기본적인 함수 메서드 바인딩 구조를 다음처럼 이해할 수 있다.

\begin{dhkdiagram}[in=2]
node a (1,1) :: Function Object + Instance Object
node b (1,3) :: Bound Method

edge a -> b :: __get__
\end{dhkdiagram}

여기서 주의할 점은 Python이 모든 속성에 대해 무조건 바운드 메서드를 생성하는 것은 아니라는 것이다.

클래스에 저장된 일반적인 함수 객체가 디스크립터 프로토콜을 제공하기 때문에 이런 동작이 가능한 것이다.

또한 인스턴스에 함수를 직접 저장한 경우에는 동작이 달라진다.

\begin{dhkcode}
def greet(self):
    print("Hello")

class Person:
    pass

p = Person()
p.greet = greet
\end{dhkcode}

이제 다음 코드를 실행하면 어떻게 될까?

\begin{dhkcode}
p.greet()
\end{dhkcode}

이 경우에는 일반적으로 \zsh{TypeError}가 발생한다.

왜냐하면 \zsh{greet} 함수 객체가 인스턴스의 \zsh{__dict__}에 직접 저장되어 있고, 인스턴스 속성으로 꺼낸 함수에는 클래스에 정의된 함수처럼 자동 바인딩이 적용되지 않기 때문이다.

따라서 다음처럼 직접 인스턴스를 전달해야 한다.

\begin{dhkcode}
p.greet(p)
\end{dhkcode}

이 차이는 매우 중요하다.

함수 객체가 어디에 저장되어 있고 어떤 속성 탐색 경로를 통해 접근되는지에 따라 메서드 바인딩 여부가 달라질 수 있다는 뜻이다.

디스크립터 프로토콜 자체는 이후 심화 학습에서 더 자세히 다룰 수 있지만, 지금은 일반적인 메서드의 바인딩 원리를 이해하는 데 집중하면 된다.

\subsubsection{\zsh{self}와 인수 전달 - 매개변수는 어떻게 바인딩되는가?}

이번에는 인수가 여러 개인 메서드를 살펴보자.

\begin{dhkcode}
class Calculator:
    def add(self, a, b):
        return a + b
    
calc = Calculator()

result = calc.add(10,20)

print(result)
\end{dhkcode}

실행 결과:

\begin{dhkzsh}
30
\end{dhkzsh}

여기서 \zsh{add()}는 세 개의 매개변수를 정의했다.

\begin{dhkcode}
def add(self, a, b):
\end{dhkcode}

하지만 호출할 때는 두 개의 인수만 전달했다.

\begin{dhkcode}
calc.add(10,20)
\end{dhkcode}

왜 가능할까?

\zsh{calc.add}가 이미 \zsh{calc} 인스턴스에 바인딩된 메서드이기 때문이다.

개념적으로 다음과 같이 대응한다.

\begin{dhkcode}
calc.add(10,20)
\end{dhkcode}

\begin{dhkcode}
Calculator.add(calc, 10, 20)
\end{dhkcode}

따라서 매개변수 바인딩은 다음과 같이 이루어진다.

\begin{dhktable}
|매개변수|전달되는 객체|
|:---|:---|
|self|calc 인스턴스|
|a|정수 객체 10|
|b|정수 객체 20|
\end{dhktable}

이 구조는 일반 함수의 인수 전달 규칙과 동일하다.

다만 바운드 메서드에서는 연결된 인스턴스가 첫 번째 인수로 제공된다는 차이가 있을 뿐이다.

\begin{dhknote}
\begin{dhksmry}[centered]
self도 일반적인 매개변수다
\end{dhksmry}
\end{dhknote}

따라서 다음 코드도 가능하다.

\begin{dhkcode}
Calculator.add(calc, 10, 20)
\end{dhkcode}

이때는 클래스에서 함수 객체에 직접 접근했으므로 \zsh{calc}를 명시적으로 전달해야 한다.

반대로 다음 코드는 오류가 발생한다.

\begin{dhkcode}
Calculator.add(10,20)
\end{dhkcode}

\zsh{self, a, b}에 필요한 인수 세 개 중 두 개만 전달했기 때문이다.

이 예제에서는 \zsh{self}라는 이름 자체가 특별한 기능을 갖는 것이 아니라, 인스턴스 메서드 바인딩 규칙과 함수의 일반적인 인수 전달 규칙이 결합되어 동작한다는 점을 기억해야 한다.

\subsubsection{인스턴스 메서드는 클래스에 저장되는가, 인스턴스에 저장되는가?}

다음 코드를 보자.

\begin{dhkcode}
class Person:
    def introduce(self):
        print("Hello")


p1 = Person()
p2 = Person()
\end{dhkcode}

여기서 \zsh{introduce} 함수 객체는 어디에 저장되어 있을까?

일반적으로 클래스의 네임스페이스에 저장된다.

\begin{dhkcode}
print("introduce" in Person.__dict__)
print("introduce" in p1.__dict__)
print("introduce" in p2.__dict__)
\end{dhkcode}

실행 결과:

\begin{dhkzsh}
True
False
False
\end{dhkzsh}

즉, \zsh{p1}과 \zsh{p2}의 인스턴스 딕셔너리에는 \zsh{introduce} 함수가 직접 저장되어 있지 않다.

하지만 다음 코드는 모두 가능하다.

\begin{dhkcode}
p1.introduce()
p2.introduce()
\end{dhkcode}

왜냐하면 Python이 인스턴스에서 속성을 찾을 때 클래스에서도 속성을 탐색하고, 클래스에 정의된 함수를 인스턴스에 바인딩해 제공하기 때문이다.

다음 코드를 확인해 보자.

\begin{dhkcode}
m1 = p1.introduce
m2 = p2.introduce

print(m1.__func__ is m2.__func__)
print(m1.__self__ is m2.__self__)
\end{dhkcode}

실행 결과:

\begin{dhkzsh}
True
False
\end{dhkzsh}

두 바운드 메서드는 같은 함수 객체를 참조하지만, 서로 다른 인스턴스에 연결되어 있다.

\begin{dhkdiagram}[in=6]
node per (1,2) :: Person
node int (3,2) :: introduce 함수 객체
node f1 (4,1) :: __func__
node f2 (4,3) :: __func__
node m1 (6,1) :: m1
node m2 (6,3) :: m2
node s1 (8,1) :: __self__
node s2 (8,3) :: __self__
node p1 (10,1) :: p1
node p2 (10,3) :: p2

edge per -> int
edge f1 -> int :: VH
edge f2 -> int :: VH
edge m1 -> f1
edge m2 -> f2
edge m1 -> s1
edge m2 -> s2
edge s1 -> p1
edge s2 -> p2
\end{dhkdiagram}

이 구조 덕분에 인스턴스마다 동일한 메서드 코드를 복사해서 저장할 필요가 없다.

하나의 함수 정의를 여러 인스턴스가 각자의 상태와 결합하여 사용할 수 있다.

\subsubsection{메서드 안에서 다른 메서드를 호출할 때도 \zsh{self}가 필요할까?}

다음 코드를 살펴보자.

\begin{dhkcode}
class Calculator:
    def add(self, a, b):
        return a + b

    def calculate(self, a, b):
        return self.add(a,b)


calc = Calculator()

print(calc.calculate(10,20))
\end{dhkcode}

실행 결과:

\begin{dhkzsh}
30
\end{dhkzsh}

여기서 \zsh{calculate()} 내부에서는 다음 코드를 사용했다.

\begin{dhkcode}
self.add(a,b)
\end{dhkcode}

왜 단순히 \zsh{add(a,b)}라고 작성하지 않을까?

Python의 일반적인 이름 탐색 규칙에서는 메서드 내부의 단순 이름 \zsh{add}를 자동으로 클래스 속성에서 찾아주지 않기 때문이다.

다음처럼 작성하면

\begin{dhkcode}
class Calculator:
    def add(self, a, b):
        return a + b

    def calculate(self, a, b):
        return add(a, b)
\end{dhkcode}

별도의 \zsh{add} 이름이 지역 또는 전역 등 일반적인 이름 탐색 경로에 정의되어 있지 않다면 \zsh{NameError}가 발생한다.

클래스에 \zsh{add}라는 메서드가 정의되어 있다는 사실만으로 메서드 본문에서 \zsh{add}라는 단순 이름이 자동으로 연결되는 것은 아니다.

따라서 다음처럼 작성해야 한다.

\begin{dhkcode}
self.add(a,b)
\end{dhkcode}

여기서는 \zsh{self}가 참조하는 인스턴스를 통해 \zsh{add} 속성을 탐색하고 바운드 메서드를 호출하는 것이다.

이 내용은 이전에 공부한 LEGB 규칙과 클래스 네임스페이스의 차이를 연결해 주는 중요한 사례다.

\subsubsection{PyTorch와의 연결 - \zsh{forward(self, x)}의 의미}

이제 PyTorch 코드에 적용해 보자.

\begin{dhkcode}
import torch
import torch.nn as nn

class MyModel(nn.Module):
    def __init__(self):
        super().__init__()
        self.linear = nn.Linear(10,5)

    def forward(self, x):
        return self.linear(x)

model = MyModel()
x = torch.randn(2, 10)

output = model(x)
\end{dhkcode}

여기서 다음 메서드에 집중해 보자.

\begin{dhkcode}
def forward(self, x):
    return self.linear(x)
\end{dhkcode}

\zsh{forward()}는 \zsh{MyModel} 클래스에 정의된 인스턴스 메서드다.

따라서 다음처럼 직접 호출할 수도 있다.

\begin{dhkcode}
output = model.forward(x)
\end{dhkcode}

이 호출은 개념적으로 다음과 대응한다.

\begin{dhkcode}
output = MyModel.forward(model, x)
\end{dhkcode}

여기서 매개변수 바인딩은 다음과 같다.

\begin{dhktable}
|매개변수|전달되는 객체|
|:---|:---|
|self|model 인스턴스|
|x|입력 Tensor 객체|
\end{dhktable}

따라서 \zsh{forward()} 내부의 다음 표현식에서

\begin{dhkcode}
self.linear(x)
\end{dhkcode}

\zsh{self}는 \zsh{model} 인스턴스를 참조한다.

그리고 \zsh{self.linear}는 해당 모델에 등록된 \zsh{nn.Linear} 하위 모듈을 찾아 반환한다.

그다음 \zsh{nn.Linear} 인스턴스를 호출해 연산을 수행하는 것이다.

\begin{dhknote}
\begin{dhksmry}[centered]
그렇다면 model(x)는 왜 가능한가?
\end{dhksmry}
\end{dhknote}

\begin{dhkcode}
output = model(x)
\end{dhkcode}

이 코드는 \zsh{forward()}를 직접 호출하는 것과는 다른 경로를 거친다.

\zsh{nn.Module}은 인스턴스를 함수처럼 호출할 수 있도록 \zsh{__call__} 메커니즘을 제공한다.

따라서 개념적인 실행 흐름은 다음과 같다.

\begin{dhksmry}[centered]
model(x)
\da
nn.Module의 __call__ 메커니즘
\da
forward(x) 호출
\da
self.linear(x)
\da
nn.Linear 연산
\da
출력 Tensor 반환
\end{dhksmry}

실제 PyTorch의 \zsh{__call__} 메커니즘은 Hook 처리 등 추가적인 기능을 포함하므로 단순히 \zsh{forward()}만 호출하는 구조는 아니다.

따라서 일반적으로는 다음 방식을 사용해야 한다.

\begin{dhkcode}
output = model(x)
\end{dhkcode}

이 구조를 정확히 이해하려면 앞으로 학습할 두 개념이 더 필요하다.

\begin{dhksmry}
    - 상속과 super() - nn.Module이 제공하는 기능을 어떻게 물려받는가?

    - 특수 메서드 - __call__()이 왜 인스턴스 호출에 사용되는가?
\end{dhksmry}

그리고 최종적으로 PyTorch \zsh{nn.Module}의 객체 모델에서 전체 구조를 연결할 것이다.

\subsubsection{핵심 정리}

이번 학습에서는 메서드와 \zsh{self}의 관계를 학습했다.

\begin{dhknote}
\begin{dhksmry}[centered]
클래스에 정의된 일반적인 인스턴스 메서드는 함수 객체로 저장된다.
\end{dhksmry}
\end{dhknote}

\begin{dhkcode}
class Person:
    def introduce(self):
        print("Hello")
\end{dhkcode}

\zsh{introduce}는 클래스 네임스페이스에 저장된 함수 객체다.

\begin{dhknote}
\begin{dhksmry}[centered]
인스턴스를 통해 메서드에 접근하면 바운드 메서드를 얻는다.
\end{dhksmry}
\end{dhknote}

\begin{dhkcode}
p.introduce
\end{dhkcode}

일반적인 경우 Python은 함수 객체와 인스턴스를 연결한 바운드 메서드를 제공한다.

\begin{dhknote}
\begin{dhksmry}[centered]
self는 Python의 예약어가 아니다.
\end{dhksmry}
\end{dhknote}

인스턴스를 전달받는 첫 번째 매개변수의 관례적인 이름이다.

\begin{dhknote}
\begin{dhksmry}[centered]
다음 두 호출은 일반적인 인스턴스 메서드에서 같은 효과를 가진다.
\end{dhksmry}
\end{dhknote}

\begin{dhkcode}
p.introduce()

Person.introduce(p)
\end{dhkcode}

단, 클래스의 함수가 재정의되거나 인스턴스 속성에 의해 가려지는 등의 경우에는 항상 같은 동작이라고 단정할 수는 없다.

\begin{dhknote}
\begin{dhksmry}[centered]
바운드 메서드는 원래 함수와 연결된 인스턴스를 참조한다.
\end{dhksmry}
\end{dhknote}

\begin{dhkcode}
method.__func__
method.__self__
\end{dhkcode}

\zsh{__func__}는 원래 함수 객체, \zsh{__self__}는 연결된 인스턴스 객체를 참조한다.

\begin{dhknote}
\begin{dhksmry}[centered]
메서드 바인딩은 디스크립터 프로토콜과 연결된다.
\end{dhksmry}
\end{dhknote}

일반적인 함수 객체는 \zsh{__get__()}을 제공하며, 이를 통해 바운드 메서드가 만들어질 수 있다.

\begin{dhknote}
\begin{dhksmry}[centered]
인스턴스 메서드는 일반적으로 클래스에 저장된다.
\end{dhksmry}
\end{dhknote}

인스턴스마다 메서드 함수 객체가 복사되는 것은 아니다.

\begin{dhknote}
\begin{dhksmry}[centered]
메서드 내부에서 다른 인스턴스 메서드를 호출할 때는 

self를 사용한다.
\end{dhksmry}
\end{dhknote}

\begin{dhkcode}
self.add(a,b)
\end{dhkcode}

이것은 인스턴스의 속성 탐색과 메서드 바인딩을 이용하는 표현식이다.

\begin{dhknote}
\begin{dhksmry}[centered]
PyTorch의 forward(self, x)도 인스턴스 메서드다.
\end{dhksmry}
\end{dhknote}

\zsh{self}는 모델 인스턴스를 참조하고, \zsh{x}는 입력 Tensor를 참조한다.

\subsubsection{2.5에서 반드시 이해해야 할 결론}

이번 학습의 핵심 구조는 다음과 같다.

\begin{dhksmry}[centered]
Class Namespace의 Function Object
+
Instance Object
\da
Descriptor Binding
\da
Bound Method
\da
호출
\da
Function Execution
\end{dhksmry}

즉, \zsh{self}가 자동으로 전달되는 것처럼 보이는 이유는 \zsh{self}라는 이름이 특별해서가 아니라, Python이 클래스의 함수 객체를 인스턴스에 바인딩한 메서드로 제공하기 때문이다.

이 개념은 함수 객체, 속성 탐색, 네임스페이스, 인수 전달이라는 이전 학습 내용이 하나로 연결되는 지점이기도 하다.

다음 2.6에서는 클래스 객체를 호출했을 때 인스턴스가 어떻게 생성되고 초기화되는지, \zsh{self}는 언제 전달되는지, \zsh{__new__()}와 \zsh{__init__()}의 역할이 왜 분리되어 있는지 살펴볼 것이다.
